\documentclass{fiche}

\surtitre{Fiche 2 · Nombres et calculs}
\titrefiche{Fractions et pourcentages}
\accroche{La fraction change de sens entre le CM2 et la 5\ieme{} : elle devient un
\emph{nombre} à part entière, un \emph{quotient}, et un \emph{opérateur}. Comprendre
ces trois visages évite la plupart des erreurs.}
\niveaux{6\ieme{} · 5\ieme{}}
\priorite{3}
\pourquoi{Passage obligé vers la proportionnalité, les pourcentages et tout le calcul de 4\ieme{}.}
\duree{30 min}
\domaine{Nombres, calcul et résolution de problèmes}

\begin{document}
\entetefiche

\begin{automatismes}[2]
  \auto[6\ieme]{$\tfrac14=0{,}25$ \quad $\tfrac12=0{,}5$ \quad $\tfrac34=0{,}75$}
  \auto[6\ieme]{$\tfrac32=1{,}5$ \quad $\tfrac42=2$ \quad $\tfrac52=2{,}5$}
  \auto[6\ieme]{$\tfrac12+\tfrac12$ \quad $1-\tfrac14$ \quad $\tfrac34-\tfrac14$}
  \auto[6\ieme]{$\tfrac23$ de 12 œufs \quad $\tfrac34$ de 10 m}
  \auto[5\ieme]{$3\times\ldots=7$ donc le nombre cherché est $\tfrac73$}
  \auto[5\ieme]{Reconnaître des fractions égales : $\tfrac23=\tfrac{10}{15}$}
  \auto[5\ieme]{$\tfrac{17}{5}=3+\tfrac25$}
  \auto[5\ieme]{$\tfrac13$ de 18 \quad $\tfrac14$ de 12}
  \auto[5\ieme]{Prendre 1 \%, 10 \% ou 50 \% d'un nombre}
  \auto[5\ieme]{$1{,}2=\tfrac{12}{10}=\tfrac65=1+\tfrac15=120\,\%$}
\end{automatismes}

\section{Les trois sens d'une fraction}

\begin{definition}[Le quotient — le sens nouveau en 6\ieme{}]
Pour $b\neq 0$, la fraction $\dfrac{a}{b}$ est le nombre qui, multiplié par $b$,
donne $a$ :
\[\frac{a}{b}\times b = a\]
Autrement dit, $\dfrac{a}{b}$ est le \motcle{résultat exact} de la division de $a$ par $b$.
\end{definition}

\begin{tableaufiche}{@{}G{30mm} Y Y@{}}
\ligneentete \entetecell{Sens} & \entetecell{Ce qu'on se dit} & \entetecell{Exemple}\\
\textbf{Partage} & « 3 parts sur 4 » & j'ai mangé $\tfrac34$ de la tarte\\
\textbf{Quotient} & « le quart de 3 » & 3 barres partagées entre 4 : $\tfrac34$ de barre chacun\\
\textbf{Opérateur} & « prendre les $\tfrac34$ de… » & $\tfrac34\times 20 = 15$\\
\end{tableaufiche}

\begin{retenir}[Le point qui débloque tout]
$\dfrac34$ de 1 \quad et \quad le quart de 3 \quad \textbf{sont le même nombre}.
C'est exactement ce que le programme de 6\ieme{} demande de comprendre.
\end{retenir}

\section{Placer une fraction sur une droite graduée}

Une fraction est un nombre : elle a donc une place sur la droite graduée.
Pour $\tfrac54$, on partage chaque unité en 4, puis on compte 5 graduations.

\begin{center}
\droitegraduee[110]{0}{3}{1}{1.25/$\frac{5}{4}$, 2.5/$\frac{5}{2}$}{}
\end{center}

\begin{retenir}
Une fraction peut être un \textbf{entier} ($\tfrac{8}{4}=2$), un \textbf{décimal}
($\tfrac34=0{,}75$), ou \textbf{ni l'un ni l'autre} ($\tfrac13=0{,}333\ldots$).
Savoir distinguer ces trois cas est un attendu de 6\ieme{}.
\end{retenir}

\section{Fractions égales et simplification}

On ne change pas un quotient en multipliant — ou en divisant — le numérateur
\emph{et} le dénominateur par un même nombre non nul :
\[\frac{a}{b}=\frac{a\times k}{b\times k}=\frac{a\div k}{b\div k}\]

\begin{methode}[Simplifier]
Chercher un diviseur commun, puis recommencer tant que c'est possible.
\[\frac{42}{56}\;\overset{\div 2}{=}\;\frac{21}{28}\;\overset{\div 7}{=}\;\frac{3}{4}\]
Les critères de divisibilité (fiche 1) servent exactement à ça.
\end{methode}

\section{Comparer des fractions}

\begin{tableaufiche}{@{}G{42mm} Y Y@{}}
\ligneentete \entetecell{Situation} & \entetecell{Méthode} & \entetecell{Exemple}\\
Même dénominateur & comparer les numérateurs & $\tfrac{2}{7}<\tfrac{5}{7}$\\
Même numérateur & le plus grand dénominateur donne le plus \emph{petit} nombre & $\tfrac{3}{4}>\tfrac{3}{8}$\\
Rien de commun & mettre au même dénominateur & $\tfrac56=\tfrac{15}{18}$ et $\tfrac79=\tfrac{14}{18}$, donc $\tfrac56>\tfrac79$\\
\end{tableaufiche}

\begin{piege}
$\tfrac18$ est plus \emph{petit} que $\tfrac14$ : plus on partage en parts nombreuses,
plus chaque part est petite. C'est contre-intuitif parce que $8>4$.
\end{piege}

\section{Additionner et soustraire}

\subsection{Même dénominateur \normalfont(6\ieme{})}
On additionne les numérateurs, on garde le dénominateur :
\[\frac{2}{5}+\frac{1}{5}=\frac{3}{5}\]

\subsection{\niv{5\ieme} Dénominateurs quelconques}

\begin{methode}
\begin{enumerate}[leftmargin=6mm,label=\textbf{\arabic*.}]
  \item Trouver un dénominateur commun (souvent, le plus grand suffit).
  \item Écrire chaque fraction avec ce dénominateur.
  \item Additionner les numérateurs.
\end{enumerate}
\vspace{0.8mm}
\[\frac34+\frac58 \;=\; \frac{6}{8}+\frac58 \;=\; \frac{11}{8}\]
\end{methode}

\begin{piege}
On n'additionne \textbf{jamais} les dénominateurs :
$\tfrac12+\tfrac13$ ne fait pas $\tfrac25$. Vérification : $\tfrac12+\tfrac13$ est
plus grand que $\tfrac12$, alors que $\tfrac25$ est plus petit.
La bonne réponse est $\tfrac36+\tfrac26=\tfrac56$.
\end{piege}

\section{Fraction d'une quantité}

Prendre $\tfrac34$ de 20, c'est multiplier : $\tfrac34\times 20=15$.
Deux chemins mènent au résultat, et les deux sont acceptés :

\begin{tableaufiche}{@{}G{46mm} Y@{}}
\ligneentete \entetecell{Chemin} & \entetecell{Calcul}\\
Diviser puis multiplier & $20\div 4=5$, puis $5\times 3=15$\\
Multiplier puis diviser & $20\times 3=60$, puis $60\div 4=15$\\
\end{tableaufiche}

Le schéma en barres, explicitement au programme de 6\ieme{}, rend la situation visible :

\begin{center}
\schemabarres{$\tfrac34$ = 15/3, reste = 5/1}{4}
\end{center}

\section{Les pourcentages}

\begin{definition}
Un pourcentage est une fraction de dénominateur 100 :
\[25\,\% = \frac{25}{100} = \frac14 = 0{,}25\]
\end{definition}

\begin{tableaufiche}{@{}G{50mm} Y Y@{}}
\ligneentete \entetecell{Question} & \entetecell{On fait} & \entetecell{Exemple}\\
Appliquer un pourcentage & multiplier & 15 \% de 240 : $\tfrac{15}{100}\times 240=36$\\
Calculer une proportion & diviser, puis $\times 100$ & 6 élèves sur 24 : $\tfrac{6}{24}=0{,}25=25\,\%$\\
Faire une réduction & retrancher & $-25\,\%$ sur 80 € : $80-20=60$ €\\
\end{tableaufiche}

\begin{retenir}[Les trois raccourcis à avoir en tête]
\[10\,\% = \text{diviser par }10 \qquad
  50\,\% = \text{la moitié} \qquad
  1\,\% = \text{diviser par }100\]
Le reste s'en déduit : 30 \% = trois fois 10 \% ; 5 \% = la moitié de 10 \%.
\end{retenir}

\begin{prolongement}
Le mot « fraction » vient d'un temps où on parlait de \emph{nombre rompu},
\emph{cassé} ou \emph{coupé}. Le programme invite à explorer les contextes
historiques — impôt, héritage, cadastre — qui ont rendu ces nombres nécessaires,
bien avant l'invention de la barre de fraction.
\end{prolongement}

\begin{videos}
  \video{mAsU47aPDXk}{Additionner des fractions}{Paul Olivier}{3:54}{760 000}
  \video{_TcFaeFb6sI}{Calculer un pourcentage — 5\ieme}{Yvan Monka}{2:58}{149 000}
  \video{YaLkcA4cIp4}{Encadrer une fraction entre deux entiers}{Maître Lucas}{4:11}{117 000}
\end{videos}

\end{document}
